\cvsection{Expériences professionnelles}

\begin{cventries}
  \cventry
    {Stagiaire ingénieur de recherche en IA}
    {INRIA}
    {Bordeaux, France}
    {Mars 2026 -- Septembre 2026}
    {
      \begin{cvitems}
        \item {Traduction d’un besoin des chercheurs en une application conçue et développée en autonomie pour faciliter la recherche d’articles scientifiques.}
        \item {Développement d’un agent IA multi-étapes sans framework d’orchestration avec l’API OpenAI, conception de prompts adaptés au cas d’usage et gestion des appels d’outils, de la mémoire et du contexte.}
        \item {Intégration d’un pipeline RAG avec chunking, embeddings et recherche sémantique dans ChromaDB. Automatisation de la collecte de métadonnées d’articles en ligne.}
        \item {Développement d’API REST Python/FastAPI documentées avec OpenAPI, stockage MongoDB et interface React/Next.js où l’utilisateur peut choisir quels messages transmettre au modèle.}
        \item {Définition d’un protocole d’évaluation de LLM, avec mesure du taux de réussite de la recherche et de la latence pour comparer les modèles et leurs limites.}
        \item {Évaluation d’OpenClaw dans un environnement Docker isolé pour étudier les risques associés aux actions d’un agent ayant accès à des outils et à des fichiers.}
        \item {Déploiement de LLM avec vLLM sur un cluster HPC équipé de GPU NVIDIA, conteneurisation de composants avec Docker et mise en place d’une CI/CD sur GitHub.}
      \end{cvitems}
    }

  \cventry
    {Stagiaire ingénieur de recherche en apprentissage par renforcement}
    {Université Karunya}
    {Coimbatore, Inde}
    {Juin 2025 -- Septembre 2025}
    {
      \begin{cvitems}
        \item {Développement d’un modèle d’apprentissage par renforcement pour la navigation autonome d’un robot hospitalier, puis déploiement et tests sur Raspberry Pi.}
      \end{cvitems}
    }
\end{cventries}
